\documentclass[aspectratio=169,10pt]{beamer}
\usepackage{fontspec}
\usepackage{polyglossia}
\setdefaultlanguage{russian}
\setotherlanguage{english}
\setmainfont{DejaVu Sans}
\setsansfont{DejaVu Sans}
\setmonofont{DejaVu Sans Mono}
\definecolor{teal}{HTML}{0A685B}
\definecolor{dark}{HTML}{153B35}
\definecolor{pale}{HTML}{E8F4EE}
\definecolor{orange}{HTML}{A95A33}
\setbeamercolor{frametitle}{fg=dark}
\setbeamercolor{title}{fg=dark}
\setbeamercolor{structure}{fg=teal}
\setbeamercolor{block title}{bg=teal,fg=white}
\setbeamercolor{block body}{bg=pale,fg=dark}
\setbeamercolor{block title alerted}{bg=orange,fg=white}
\setbeamercolor{block body alerted}{bg=orange!9,fg=dark}
\setbeamertemplate{navigation symbols}{}
\setbeamertemplate{footline}{\hfill\scriptsize\insertframenumber\, / \,\inserttotalframenumber\hspace{12pt}\vspace{5pt}}
\setbeamertemplate{itemize item}{\raisebox{.2ex}{\tiny$\blacktriangleright$}}
\setbeamerfont{frametitle}{size=\Large,series=\bfseries}
\setbeamerfont{normal text}{size=\normalsize}
\setlength{\leftmargini}{1.3em}
\begin{document}
\begin{frame}[plain]
\begin{beamercolorbox}[wd=\paperwidth,ht=\paperheight,sep=1.4cm,center]{title}
\vspace*{1.0cm}{\large Данные как основа ИИ · семестр 2}\par\vspace{1cm}
{\fontsize{30}{38}\selectfont\bfseries Одни данные — разные признаки}\par\vspace{0.6cm}
{\large Почему представление зависит от вопроса}\par\vfill{\small Кейс «Умный кампус» · университетские пространства}
\end{beamercolorbox}\end{frame}
\begin{frame}[t]{Одна таблица, два заказчика}
\begin{itemize}
\item Управляющий хочет понятный отчёт по аудиториям.
\vspace{4pt}
\item Инженер хочет сравнивать похожие состояния помещений.
\vspace{4pt}
\item Одна и та же исходная запись может иметь разные полезные представления.
\vspace{4pt}
\end{itemize}
\vfill\begin{alertblock}{Вопрос аудитории}Должны ли выходные столбцы быть одинаковыми?\end{alertblock}
\end{frame}
\begin{frame}[t]{Данные и смысловые типы}
\begin{itemize}
\item record\_id — идентификатор, не количество.
\vspace{4pt}
\item room\_type — номинальная категория.
\vspace{4pt}
\item co2\_ppm — непрерывное измерение с единицами.
\vspace{4pt}
\item observed\_at — время, из которого получаются календарные признаки.
\vspace{4pt}
\end{itemize}
\end{frame}
\begin{frame}[t]{Сравнивать занятость корректно}
\begin{itemize}
\item 18 человек из 20 мест: доля заполнения 0,9.
\vspace{4pt}
\item 18 человек из 60 мест: доля заполнения 0,3.
\vspace{4pt}
\item Смысловой признак важнее произвольного числового кода комнаты.
\vspace{4pt}
\end{itemize}
\vfill\begin{alertblock}{Вопрос аудитории}Что делать, если вместимость отсутствует?\end{alertblock}
\end{frame}
\begin{frame}[t]{Признак должен знать свою формулу}
\begin{itemize}
\item day = календарная дата измерения.
\vspace{4pt}
\item hour = час измерения.
\vspace{4pt}
\item occupancy\_ratio = people\_count / capacity.
\vspace{4pt}
\item quality\_flag сохраняет ограничения исходного наблюдения.
\vspace{4pt}
\end{itemize}
\end{frame}
\begin{frame}[t]{Порог — не закон природы}
\begin{itemize}
\item Отчёт может использовать учебный индикатор «значение выше X».
\vspace{4pt}
\item X должен быть указан и обоснован.
\vspace{4pt}
\item Учебные пороги не заменяют действующие санитарные нормы и калибровку.
\vspace{4pt}
\end{itemize}
\end{frame}
\begin{frame}[t]{Номинальная категория не является шкалой}
\begin{itemize}
\item Лекционная=1, семинарская=2, лабораторная=3 — произвольные коды.
\vspace{4pt}
\item Алгоритм расстояния увидит несуществующий порядок.
\vspace{4pt}
\item Для номинальных категорий подойдёт индикаторное представление.
\vspace{4pt}
\end{itemize}
\vfill\begin{alertblock}{Вопрос аудитории}Почему лабораторная «дальше» от лекционной на один шаг?\end{alertblock}
\end{frame}
\begin{frame}[t]{Когда масштаб мешает}
\begin{itemize}
\item Вектор содержит CO₂ около 900, температуру 22 и заполнение 0,7.
\vspace{4pt}
\item В прямом расстоянии вклад CO₂ может доминировать.
\vspace{4pt}
\item При масштабировании надо сохранить смысл и зафиксировать параметры.
\vspace{4pt}
\end{itemize}
\end{frame}
\begin{frame}[t]{Min–max без магии}
\begin{itemize}
\item Формула: новое = (x−min) / (max−min).
\vspace{4pt}
\item На новой партии не нужно пересчитывать min и max заново, если требуются сопоставимые значения.
\vspace{4pt}
\item Результат за пределами 0–1 возможен для нового экстремума.
\vspace{4pt}
\end{itemize}
\end{frame}
\begin{frame}[t]{Стандартизация и выбросы}
\begin{itemize}
\item Z-оценка — отклонение от среднего в единицах стандартного отклонения.
\vspace{4pt}
\item Экстремальные наблюдения влияют на среднее и разброс.
\vspace{4pt}
\item Масштабирование ID и произвольных категорий не придаёт им смысла.
\vspace{4pt}
\end{itemize}
\end{frame}
\begin{frame}[t]{Представление А: читаемый отчёт}
\begin{itemize}
\item room\_id, корпус, дата, час, CO₂ ppm, °C, заполнение.
\vspace{4pt}
\item Видны статусы качества и пропуски.
\vspace{4pt}
\item Названия категорий сохраняются для человека.
\vspace{4pt}
\end{itemize}
\end{frame}
\begin{frame}[t]{Представление Б: численная схема}
\begin{itemize}
\item Численные поля согласованы по шкале.
\vspace{4pt}
\item Номинальные категории закодированы без ложного порядка.
\vspace{4pt}
\item ID остаётся для идентификации, но не участвует в расстоянии.
\vspace{4pt}
\end{itemize}
\end{frame}
\begin{frame}[t]{Из сообщения — в структурированный признак}
\begin{itemize}
\item «В аудитории душно» -> тема «воздух».
\vspace{4pt}
\item «Датчик не отвечает» -> тема «датчик» или нужна проверка.
\vspace{4pt}
\item LLM предлагает метку, проверка идёт по закрытому словарю и ID.
\vspace{4pt}
\end{itemize}
\vfill\begin{alertblock}{Вопрос аудитории}Как проверить, что модель не потеряла ни одного сообщения?\end{alertblock}
\end{frame}
\begin{frame}[t]{Что нельзя принимать от ИИ}
\begin{itemize}
\item «Нормализуем номер аудитории как число».
\vspace{4pt}
\item «Всем категориям присвоим 1–2–3».
\vspace{4pt}
\item «Всё выше 1.0 по заполнению обрежем».
\vspace{4pt}
\item «На каждом файле пересчитаем новую шкалу».
\vspace{4pt}
\end{itemize}
\vfill\begin{alertblock}{Вопрос аудитории}Выберите самый опасный пункт для нашей задачи и объясните.\end{alertblock}
\end{frame}
\begin{frame}[t]{Лабораторная: две версии}
\begin{itemize}
\item Результаты А и Б строятся из одной исходной таблицы.
\vspace{4pt}
\item Словарь признаков: смысл, формула, единицы, ограничения.
\vspace{4pt}
\item Пять контрольных строк + аудит рекомендаций ИИ.
\vspace{4pt}
\item Необязательное углубление — LLM-разметка сообщений.
\vspace{4pt}
\end{itemize}
\end{frame}
\begin{frame}[t]{Работоспособность новой партии}
\begin{itemize}
\item Сохранили ли мы правила и параметры?
\vspace{4pt}
\item Сохранились ли ID и число наблюдений?
\vspace{4pt}
\item Как вести себя при неизвестной категории или пустой вместимости?
\vspace{4pt}
\end{itemize}
\end{frame}
\begin{frame}[t]{Выходной билет}
\begin{itemize}
\item Назовите преобразование, которое полезно для отчёта, но не нужно для сравнения.
\vspace{4pt}
\item Назовите операцию, способную создать ложное числовое расстояние.
\vspace{4pt}
\item Что вы должны проверить, прежде чем доверять признаку, извлечённому LLM?
\vspace{4pt}
\end{itemize}
\end{frame}
\end{document}
